\chapter*{Notation and Conventions}
\phantomsection
\label{sec:i-notation-and-conventions}
\addcontentsline{toc}{chapter}{Notation and Conventions}

\section*{Standing conventions}

Unless a section explicitly states otherwise, Volume~I uses the following
conventions.

\begin{enumerate}
\item Vector spaces in the algebraic chapters are finite-dimensional whenever
a dimension-dependent theorem or construction is invoked. Statements that do
not require finite dimension remain valid without that restriction.

\item The scalar field is denoted by \(\mathbb F\), usually
\(\mathbb R\) or \(\mathbb C\). When the distinction matters, the field is
stated explicitly.

\item Complex inner products are linear in the first argument and
conjugate-linear in the second:
\[
\langle \alpha x,y\rangle
 = \alpha\langle x,y\rangle,
\qquad
\langle x,\alpha y\rangle
 = \overline{\alpha}\langle x,y\rangle.
\]
Consequently,
\[
\langle x,y\rangle=\overline{\langle y,x\rangle}.
\]

\item For a complex matrix,
\[
A^*=\overline A^{\,T}
\]
denotes the conjugate transpose; over \(\mathbb R\), this reduces to
\(A^T\). For a linear operator on an inner-product space, \(T^*\) denotes
the adjoint characterized by
\[
\langle Tv,w\rangle=\langle v,T^*w\rangle.
\]

\item On \(\mathbb C^n\), the standard inner product is therefore written
\[
\langle x,y\rangle=y^*x.
\]
This ordering is consistent with linearity in the first argument.
\end{enumerate}

These conventions govern the later treatment of orthogonality, adjoints,
Hermitian and unitary matrices, normal operators, the spectral theorem, and
singular-value decomposition. A chapter repeats a convention only when the
placement of conjugation or the scalar field enters a calculation directly.
